\chapter{Atypical Chest Pain}

\section*{Definition}
Atypical chest pain is chest discomfort that does not have the classic pattern of pressure brought on by exertion and relieved by rest. The term does not mean harmless pain. People with heart disease, particularly women, older adults, and people with diabetes, may have atypical or associated symptoms.

\section*{What Happens in Your Body}
Chest symptoms can arise from the heart, lungs, aorta, esophagus, chest wall, skin, or nerves. Myocardial ischemia may cause pressure, burning, indigestion, breathlessness, nausea, fatigue, or pain in the jaw, back, or upper abdomen rather than obvious central chest pain.

\section*{Causes}
\begin{itemize}
  \item Acute coronary syndrome, stable angina, coronary vasospasm, or microvascular angina
  \item Pulmonary embolism, pneumothorax, pneumonia, or pericarditis
  \item Aortic dissection or other vascular emergency
  \item Reflux, esophageal spasm, peptic disease, or gallbladder disease
  \item Costochondritis, muscle strain, panic attack, or anxiety
\end{itemize}

\section*{When to See a Doctor}
Call emergency services for new or unexplained chest symptoms with shortness of breath, sweating, nausea, fainting, palpitations, weakness, or pain radiating to the arm, jaw, back, or upper abdomen. Do not dismiss symptoms as indigestion or anxiety without assessment when risk is significant.

\section*{Related Symptoms}
\begin{itemize}
  \item Fatigue, breathlessness, dizziness, nausea, sweating, or palpitations
  \item Pain with breathing, cough, fever, wheeze, or coughing blood
  \item Tenderness with touch or movement, heartburn, or regurgitation
\end{itemize}

\section*{Self-Care / Home Management}
Stop exertion and seek emergency advice for concerning symptoms. Do not drive yourself, take someone else’s medication, or use an antacid trial to delay evaluation. Use prescribed cardiac or respiratory rescue medicine only according to an established plan.

\section*{How Your Doctor Will Evaluate You}
\begin{itemize}
  \item Initial evaluation includes vital signs, oxygen saturation, cardiovascular and lung examination, and a 12-lead ECG.
  \item Serial high-sensitivity troponins, blood tests, and chest imaging assess cardiac, pulmonary, and systemic causes.
  \item Echocardiography, CT angiography, stress testing, or coronary angiography is selected according to risk and initial results.
\end{itemize}

\section*{Treatment Options}
Treatment may include emergency coronary reperfusion, antiplatelet therapy, anticoagulation, treatment of pulmonary disease, antibiotics, acid suppression, or musculoskeletal and anxiety care after dangerous causes are excluded.

\section*{References}
\begin{thebibliography}{99}
\bibitem{atypical_chest_pain:ref1} Gulati M, et al. 2021 AHA/ACC Guideline for the Evaluation and Diagnosis of Chest Pain. \textit{Circulation}. 2021;144:e368--e454.
\bibitem{atypical_chest_pain:ref2} Byrne RA, et al. 2023 ESC Guidelines for acute coronary syndromes. \textit{Eur Heart J}. 2023;44:3720--3826.
\end{thebibliography}
